% --- Archon physics formalization source begin ---
% archon:physics
% archon:covers PhyXMiniProblems/problem_phyx_mini_0388.lean
% archon:source-report reports/phyx_mini/problem_phyx_mini_0388.source.json

\chapter{Physics problem phyx\_mini\_0388}
\label{ch:PhyXMiniProblems_problem_phyx_mini_0388}

\paragraph{Problem source.}
\#\# Physical scenario

A pipe flowing light oil has a manometer attached.

\#\# Question

What is the absolute pressure in the pipe flow?

\#\# Answer choices

- A: A: 490kPa

- B: B: 106.4kPa

- C: C: 154kPa

- D: D: 10.2kPa

\#\# Figure caption (auxiliary; use the image as primary evidence)

The image depicts a U-tube manometer setup, designed to measure pressure differences. Here's a detailed description:

1. **U-tube Design**: 
   - The U-tube contains two different fluids: oil and water. 
   - Water is at the bottom, with oil layered above it on the left side.

2. **Dimensions and Labels**:
   - The height of the oil column is labeled as 0.3 m above the base.
   - The top of the U-tube on the right side has a height of 0.7 m from the base.
   - There's another measurement showing the height of the water column as 0.1 m.

3. **Pressure Details**: 
   - External pressure at the top of the right side of the manometer is indicated as \textbackslash{}( P\_0 = 101 \textbackslash{}, \textbackslash{}text\{kPa\} \textbackslash{}).

4. **Fluid Flow**:
   - An arrow shows the direction of fluid pressure entering the left side of the tube.

5. **Visual Representation**:
   - The oil is colored to differentiate it from the water.
   - The direction of flow and differences in fluid levels are clearly shown to convey pressure differences.

This setup suggests a measurement of the pressure difference between two points using the density difference between oil and water.

\#\# Dataset metadata

Temperature and Heat Transfer; Physical Model Grounding Reasoning, Multi-Formula Reasoning

\paragraph{Recorded answer/context.}
B: B: 106.4kPa

\paragraph{Figure/image path.}
/root/proposal\_for\_physic/science-mango/phyx\_mini\_run/phyx\_data/test\_image/388.png

\paragraph{Formalization target.}
create a compiling Lean file with sorry bodies at `PhyXMiniProblems/problem\_phyx\_mini\_0388.lean`. The Lean declarations must preserve the physical quantities, dimensions or dimensional roles, figure labels, governing-law hypotheses, and final relation expressed by this problem.
Use Mathlib/Physlib names found through LeanExplore where available. If a physics API is missing, introduce faithful local abstractions rather than scalar placeholder aliases.

\begin{theorem}[Physics formalization target]
\label{thm:physics:phyx_mini_0388:target}
\lean{PhyXMiniProblems.ProblemPhyXMini0388.problem_phyx_mini_0388}
\uses{def:physics:phyx-mini-0388:phyxminiproblems-problemphyxmini0388-pressureinpascals, def:physics:phyx-mini-0388:phyxminiproblems-problemphyxmini0388-pressureinkilopascals, def:physics:phyx-mini-0388:phyxminiproblems-problemphyxmini0388-pipemanometersetup, def:physics:phyx-mini-0388:phyxminiproblems-problemphyxmini0388-matchesprimaryfigure, def:physics:phyx-mini-0388:phyxminiproblems-problemphyxmini0388-usesexplicitnumericalcalibration, def:physics:phyx-mini-0388:phyxminiproblems-problemphyxmini0388-hasphysicalmanometerparameters, def:physics:phyx-mini-0388:phyxminiproblems-problemphyxmini0388-satisfiestwofluidhydrostaticlaws, def:physics:phyx-mini-0388:phyxminiproblems-problemphyxmini0388-isuniquematchingabsolutepressure}
The assigned autoformalize agent should translate this physics problem into Lean declarations in the covered file, with theorem and lemma proof bodies written as `by sorry`.
\end{theorem}
\begin{proof}
This is an autoformalization task, not a proof task. Produce statements that can later be proved by the physics prover without weakening the physical model.
\end{proof}

% --- Archon declaration topology begin ---
\section{Declaration topology}
\begin{definition}[mass Density Dimension]
  \label{def:physics:phyx-mini-0388:phyxminiproblems-problemphyxmini0388-massdensitydimension}
  \lean{PhyXMiniProblems.ProblemPhyXMini0388.massDensityDimension}
  The physical dimension of mass density, M L⁻³
\end{definition}
\begin{proof}
  This is the declared model definition of mass Density Dimension.
\end{proof}
\begin{definition}[acceleration Dimension]
  \label{def:physics:phyx-mini-0388:phyxminiproblems-problemphyxmini0388-accelerationdimension}
  \lean{PhyXMiniProblems.ProblemPhyXMini0388.accelerationDimension}
  The physical dimension of acceleration, L T⁻²
\end{definition}
\begin{proof}
  This is the declared model definition of acceleration Dimension.
\end{proof}
\begin{definition}[Length Quantity]
  \label{def:physics:phyx-mini-0388:phyxminiproblems-problemphyxmini0388-lengthquantity}
  \lean{PhyXMiniProblems.ProblemPhyXMini0388.LengthQuantity}
  A nonnegative physical length, independent of the unit used to read it
\end{definition}
\begin{proof}
  This is the declared model definition of Length Quantity.
\end{proof}
\begin{definition}[Mass Density Quantity]
  \label{def:physics:phyx-mini-0388:phyxminiproblems-problemphyxmini0388-massdensityquantity}
  \lean{PhyXMiniProblems.ProblemPhyXMini0388.MassDensityQuantity}
  \uses{def:physics:phyx-mini-0388:phyxminiproblems-problemphyxmini0388-massdensitydimension}
  A nonnegative physical mass density
\end{definition}
\begin{proof}
  This is the declared model definition of Mass Density Quantity.
\end{proof}
\begin{definition}[Acceleration Quantity]
  \label{def:physics:phyx-mini-0388:phyxminiproblems-problemphyxmini0388-accelerationquantity}
  \lean{PhyXMiniProblems.ProblemPhyXMini0388.AccelerationQuantity}
  \uses{def:physics:phyx-mini-0388:phyxminiproblems-problemphyxmini0388-accelerationdimension}
  A nonnegative physical acceleration magnitude
\end{definition}
\begin{proof}
  This is the declared model definition of Acceleration Quantity.
\end{proof}
\begin{definition}[Pressure Quantity]
  \label{def:physics:phyx-mini-0388:phyxminiproblems-problemphyxmini0388-pressurequantity}
  \lean{PhyXMiniProblems.ProblemPhyXMini0388.PressureQuantity}
  Absolute pressure, represented using Physlib's pressure dimension
\end{definition}
\begin{proof}
  This is the declared model definition of Pressure Quantity.
\end{proof}
\begin{definition}[length In Meters]
  \label{def:physics:phyx-mini-0388:phyxminiproblems-problemphyxmini0388-lengthinmeters}
  \lean{PhyXMiniProblems.ProblemPhyXMini0388.lengthInMeters}
  \uses{def:physics:phyx-mini-0388:phyxminiproblems-problemphyxmini0388-lengthquantity}
  Metre readout of a physical length
\end{definition}
\begin{proof}
  This is the declared model definition of length In Meters.
\end{proof}
\begin{definition}[density In Kilograms Per Cubic Meter]
  \label{def:physics:phyx-mini-0388:phyxminiproblems-problemphyxmini0388-densityinkilogramspercubicmeter}
  \lean{PhyXMiniProblems.ProblemPhyXMini0388.densityInKilogramsPerCubicMeter}
  \uses{def:physics:phyx-mini-0388:phyxminiproblems-problemphyxmini0388-massdensityquantity}
  Kilogram-per-cubic-metre readout of a physical mass density
\end{definition}
\begin{proof}
  This is the declared model definition of density In Kilograms Per Cubic Meter.
\end{proof}
\begin{definition}[acceleration In Meters Per Second Squared]
  \label{def:physics:phyx-mini-0388:phyxminiproblems-problemphyxmini0388-accelerationinmeterspersecondsquared}
  \lean{PhyXMiniProblems.ProblemPhyXMini0388.accelerationInMetersPerSecondSquared}
  \uses{def:physics:phyx-mini-0388:phyxminiproblems-problemphyxmini0388-accelerationquantity}
  Metre-per-second-squared readout of an acceleration magnitude
\end{definition}
\begin{proof}
  This is the declared model definition of acceleration In Meters Per Second Squared.
\end{proof}
\begin{definition}[pressure In Pascals]
  \label{def:physics:phyx-mini-0388:phyxminiproblems-problemphyxmini0388-pressureinpascals}
  \lean{PhyXMiniProblems.ProblemPhyXMini0388.pressureInPascals}
  \uses{def:physics:phyx-mini-0388:phyxminiproblems-problemphyxmini0388-pressurequantity}
  Pascal readout of an absolute pressure
\end{definition}
\begin{proof}
  This is the declared model definition of pressure In Pascals.
\end{proof}
\begin{definition}[pressure In Kilopascals]
  \label{def:physics:phyx-mini-0388:phyxminiproblems-problemphyxmini0388-pressureinkilopascals}
  \lean{PhyXMiniProblems.ProblemPhyXMini0388.pressureInKilopascals}
  \uses{def:physics:phyx-mini-0388:phyxminiproblems-problemphyxmini0388-pressurequantity, def:physics:phyx-mini-0388:phyxminiproblems-problemphyxmini0388-pressureinpascals}
  Kilopascal readout of an absolute pressure
\end{definition}
\begin{proof}
  This is the declared model definition of pressure In Kilopascals.
\end{proof}
\begin{definition}[Manometer Arm]
  \label{def:physics:phyx-mini-0388:phyxminiproblems-problemphyxmini0388-manometerarm}
  \lean{PhyXMiniProblems.ProblemPhyXMini0388.ManometerArm}
  The two arms of the U-tube, named by their boundary conditions
\end{definition}
\begin{proof}
  This is the declared model definition of Manometer Arm.
\end{proof}
\begin{definition}[Manometer Fluid]
  \label{def:physics:phyx-mini-0388:phyxminiproblems-problemphyxmini0388-manometerfluid}
  \lean{PhyXMiniProblems.ProblemPhyXMini0388.ManometerFluid}
  Fluids named in the problem and primary image
\end{definition}
\begin{proof}
  This is the declared model definition of Manometer Fluid.
\end{proof}
\begin{definition}[Figure Height Label]
  \label{def:physics:phyx-mini-0388:phyxminiproblems-problemphyxmini0388-figureheightlabel}
  \lean{PhyXMiniProblems.ProblemPhyXMini0388.FigureHeightLabel}
  The three vertical dimensions printed in the image
\end{definition}
\begin{proof}
  This is the declared model definition of Figure Height Label.
\end{proof}
\begin{definition}[Primary Manometer Figure]
  \label{def:physics:phyx-mini-0388:phyxminiproblems-problemphyxmini0388-primarymanometerfigure}
  \lean{PhyXMiniProblems.ProblemPhyXMini0388.PrimaryManometerFigure}
  \uses{def:physics:phyx-mini-0388:phyxminiproblems-problemphyxmini0388-lengthquantity, def:physics:phyx-mini-0388:phyxminiproblems-problemphyxmini0388-pressurequantity, def:physics:phyx-mini-0388:phyxminiproblems-problemphyxmini0388-manometerarm, def:physics:phyx-mini-0388:phyxminiproblems-problemphyxmini0388-manometerfluid, def:physics:phyx-mini-0388:phyxminiproblems-problemphyxmini0388-figureheightlabel}
  Structured transcription of image 388.png. The 0.7 m arrow extends from the oil--water interface level to the right free surface; it is not measured from the bottom datum
\end{definition}
\begin{proof}
  This is the declared model definition of Primary Manometer Figure.
\end{proof}
\begin{definition}[Pipe Manometer Setup]
  \label{def:physics:phyx-mini-0388:phyxminiproblems-problemphyxmini0388-pipemanometersetup}
  \lean{PhyXMiniProblems.ProblemPhyXMini0388.PipeManometerSetup}
  \uses{def:physics:phyx-mini-0388:phyxminiproblems-problemphyxmini0388-lengthquantity, def:physics:phyx-mini-0388:phyxminiproblems-problemphyxmini0388-massdensityquantity, def:physics:phyx-mini-0388:phyxminiproblems-problemphyxmini0388-accelerationquantity, def:physics:phyx-mini-0388:phyxminiproblems-problemphyxmini0388-pressurequantity, def:physics:phyx-mini-0388:phyxminiproblems-problemphyxmini0388-manometerfluid, def:physics:phyx-mini-0388:phyxminiproblems-problemphyxmini0388-primarymanometerfigure}
  The apparatus and its independent physical observables. The pressure at the oil--water interface is explicit so that hydrostatics can be imposed once through each fluid column. No requested pressure value is defined here
\end{definition}
\begin{proof}
  This is the declared model definition of Pipe Manometer Setup.
\end{proof}
\begin{definition}[Matches Primary Figure]
  \label{def:physics:phyx-mini-0388:phyxminiproblems-problemphyxmini0388-matchesprimaryfigure}
  \lean{PhyXMiniProblems.ProblemPhyXMini0388.MatchesPrimaryFigure}
  \uses{def:physics:phyx-mini-0388:phyxminiproblems-problemphyxmini0388-lengthinmeters, def:physics:phyx-mini-0388:phyxminiproblems-problemphyxmini0388-pressureinpascals, def:physics:phyx-mini-0388:phyxminiproblems-problemphyxmini0388-pipemanometersetup}
  Qualitative assignments and scalar readouts taken directly from the primary image. In particular, the atmosphere label is an absolute pressure. This predicate contains no pipe-pressure answer
\end{definition}
\begin{proof}
  This is the declared model definition of Matches Primary Figure.
\end{proof}
\begin{definition}[Uses Explicit Numerical Calibration]
  \label{def:physics:phyx-mini-0388:phyxminiproblems-problemphyxmini0388-usesexplicitnumericalcalibration}
  \lean{PhyXMiniProblems.ProblemPhyXMini0388.UsesExplicitNumericalCalibration}
  \uses{def:physics:phyx-mini-0388:phyxminiproblems-problemphyxmini0388-densityinkilogramspercubicmeter, def:physics:phyx-mini-0388:phyxminiproblems-problemphyxmini0388-accelerationinmeterspersecondsquared, def:physics:phyx-mini-0388:phyxminiproblems-problemphyxmini0388-pipemanometersetup}
  Explicit numerical calibration needed to recover the recorded multiple-choice answer: water has density 1000 kg/m³, light oil has density 800 kg/m³, and the exercise uses g = 10 m/s². The problem source does not state these density/gravity values, so this calibration is deliberately separate from the figure transcription and the source-determined symbolic result below
\end{definition}
\begin{proof}
  This is the declared model definition of Uses Explicit Numerical Calibration.
\end{proof}
\begin{definition}[Has Physical Manometer Parameters]
  \label{def:physics:phyx-mini-0388:phyxminiproblems-problemphyxmini0388-hasphysicalmanometerparameters}
  \lean{PhyXMiniProblems.ProblemPhyXMini0388.HasPhysicalManometerParameters}
  \uses{def:physics:phyx-mini-0388:phyxminiproblems-problemphyxmini0388-lengthinmeters, def:physics:phyx-mini-0388:phyxminiproblems-problemphyxmini0388-densityinkilogramspercubicmeter, def:physics:phyx-mini-0388:phyxminiproblems-problemphyxmini0388-accelerationinmeterspersecondsquared, def:physics:phyx-mini-0388:phyxminiproblems-problemphyxmini0388-pressureinpascals, def:physics:phyx-mini-0388:phyxminiproblems-problemphyxmini0388-pipemanometersetup}
  Positivity and geometric ordering required for the static-fluid model
\end{definition}
\begin{proof}
  This is the declared model definition of Has Physical Manometer Parameters.
\end{proof}
\begin{definition}[Satisfies Two Fluid Hydrostatic Laws]
  \label{def:physics:phyx-mini-0388:phyxminiproblems-problemphyxmini0388-satisfiestwofluidhydrostaticlaws}
  \lean{PhyXMiniProblems.ProblemPhyXMini0388.SatisfiesTwoFluidHydrostaticLaws}
  \uses{def:physics:phyx-mini-0388:phyxminiproblems-problemphyxmini0388-lengthinmeters, def:physics:phyx-mini-0388:phyxminiproblems-problemphyxmini0388-densityinkilogramspercubicmeter, def:physics:phyx-mini-0388:phyxminiproblems-problemphyxmini0388-accelerationinmeterspersecondsquared, def:physics:phyx-mini-0388:phyxminiproblems-problemphyxmini0388-pressureinpascals, def:physics:phyx-mini-0388:phyxminiproblems-problemphyxmini0388-pipemanometersetup}
  For a static fluid, pressure increases by ρ g Δh when descending. Applied from the pipe tap through the oil and from the open free surface through the water, both routes reach the same oil--water interface pressure. These are governing relations, not the solved pipe-pressure formula
\end{definition}
\begin{proof}
  This is the declared model definition of Satisfies Two Fluid Hydrostatic Laws.
\end{proof}
\begin{lemma}[pipe pressure from column balance]
  \label{lem:physics:phyx-mini-0388:phyxminiproblems-problemphyxmini0388-pipe-pressure-from-column-balance}
  \lean{PhyXMiniProblems.ProblemPhyXMini0388.pipe_pressure_from_column_balance}
  \uses{def:physics:phyx-mini-0388:phyxminiproblems-problemphyxmini0388-lengthinmeters, def:physics:phyx-mini-0388:phyxminiproblems-problemphyxmini0388-densityinkilogramspercubicmeter, def:physics:phyx-mini-0388:phyxminiproblems-problemphyxmini0388-accelerationinmeterspersecondsquared, def:physics:phyx-mini-0388:phyxminiproblems-problemphyxmini0388-pressureinpascals, def:physics:phyx-mini-0388:phyxminiproblems-problemphyxmini0388-pipemanometersetup, def:physics:phyx-mini-0388:phyxminiproblems-problemphyxmini0388-satisfiestwofluidhydrostaticlaws}
  Eliminating the common interface pressure gives the signed two-fluid manometer relation for the unknown pipe pressure
\end{lemma}
\begin{proof}
  The conclusion follows from the governing relations and figure readouts named in the statement of pipe pressure from column balance.
\end{proof}
\begin{lemma}[pipe pressure from primary figure]
  \label{lem:physics:phyx-mini-0388:phyxminiproblems-problemphyxmini0388-pipe-pressure-from-primary-figure}
  \lean{PhyXMiniProblems.ProblemPhyXMini0388.pipe_pressure_from_primary_figure}
  \uses{def:physics:phyx-mini-0388:phyxminiproblems-problemphyxmini0388-densityinkilogramspercubicmeter, def:physics:phyx-mini-0388:phyxminiproblems-problemphyxmini0388-accelerationinmeterspersecondsquared, def:physics:phyx-mini-0388:phyxminiproblems-problemphyxmini0388-pressureinpascals, def:physics:phyx-mini-0388:phyxminiproblems-problemphyxmini0388-pipemanometersetup, def:physics:phyx-mini-0388:phyxminiproblems-problemphyxmini0388-matchesprimaryfigure, def:physics:phyx-mini-0388:phyxminiproblems-problemphyxmini0388-satisfiestwofluidhydrostaticlaws}
  The strongest pressure relation determined by the primary image and the hydrostatic laws alone. It specializes the atmospheric pressure and three figure heights while leaving the omitted fluid densities and gravitational acceleration as physical parameters
\end{lemma}
\begin{proof}
  The conclusion follows from the governing relations and figure readouts named in the statement of pipe pressure from primary figure.
\end{proof}
\begin{lemma}[pipe absolute pressure in pascals]
  \label{lem:physics:phyx-mini-0388:phyxminiproblems-problemphyxmini0388-pipe-absolute-pressure-in-pascals}
  \lean{PhyXMiniProblems.ProblemPhyXMini0388.pipe_absolute_pressure_in_pascals}
  \uses{def:physics:phyx-mini-0388:phyxminiproblems-problemphyxmini0388-pressureinpascals, def:physics:phyx-mini-0388:phyxminiproblems-problemphyxmini0388-pipemanometersetup, def:physics:phyx-mini-0388:phyxminiproblems-problemphyxmini0388-matchesprimaryfigure, def:physics:phyx-mini-0388:phyxminiproblems-problemphyxmini0388-usesexplicitnumericalcalibration, def:physics:phyx-mini-0388:phyxminiproblems-problemphyxmini0388-satisfiestwofluidhydrostaticlaws}
  Substituting the image readouts and explicit calibration gives 106400 Pa
\end{lemma}
\begin{proof}
  The conclusion follows from the governing relations and figure readouts named in the statement of pipe absolute pressure in pascals.
\end{proof}
\begin{definition}[Answer Choice]
  \label{def:physics:phyx-mini-0388:phyxminiproblems-problemphyxmini0388-answerchoice}
  \lean{PhyXMiniProblems.ProblemPhyXMini0388.AnswerChoice}
  Labels of the four pressure choices printed with the problem
\end{definition}
\begin{proof}
  This is the declared model definition of Answer Choice.
\end{proof}
\begin{definition}[displayed Absolute Pressure Kilopascals]
  \label{def:physics:phyx-mini-0388:phyxminiproblems-problemphyxmini0388-displayedabsolutepressurekilopascals}
  \lean{PhyXMiniProblems.ProblemPhyXMini0388.displayedAbsolutePressureKilopascals}
  \uses{def:physics:phyx-mini-0388:phyxminiproblems-problemphyxmini0388-answerchoice}
  Absolute-pressure value in kilopascals displayed beside each choice
\end{definition}
\begin{proof}
  This is the declared model definition of displayed Absolute Pressure Kilopascals.
\end{proof}
\begin{definition}[Matches Displayed Absolute Pressure]
  \label{def:physics:phyx-mini-0388:phyxminiproblems-problemphyxmini0388-matchesdisplayedabsolutepressure}
  \lean{PhyXMiniProblems.ProblemPhyXMini0388.MatchesDisplayedAbsolutePressure}
  \uses{def:physics:phyx-mini-0388:phyxminiproblems-problemphyxmini0388-pressureinkilopascals, def:physics:phyx-mini-0388:phyxminiproblems-problemphyxmini0388-pipemanometersetup, def:physics:phyx-mini-0388:phyxminiproblems-problemphyxmini0388-answerchoice, def:physics:phyx-mini-0388:phyxminiproblems-problemphyxmini0388-displayedabsolutepressurekilopascals}
  A displayed choice agrees exactly with the modeled pressure readout
\end{definition}
\begin{proof}
  This is the declared model definition of Matches Displayed Absolute Pressure.
\end{proof}
\begin{definition}[Is Unique Matching Absolute Pressure]
  \label{def:physics:phyx-mini-0388:phyxminiproblems-problemphyxmini0388-isuniquematchingabsolutepressure}
  \lean{PhyXMiniProblems.ProblemPhyXMini0388.IsUniqueMatchingAbsolutePressure}
  \uses{def:physics:phyx-mini-0388:phyxminiproblems-problemphyxmini0388-pipemanometersetup, def:physics:phyx-mini-0388:phyxminiproblems-problemphyxmini0388-answerchoice, def:physics:phyx-mini-0388:phyxminiproblems-problemphyxmini0388-matchesdisplayedabsolutepressure}
  A choice is the unique displayed pressure matching the physical model
\end{definition}
\begin{proof}
  This is the declared model definition of Is Unique Matching Absolute Pressure.
\end{proof}
% --- Archon declaration topology end ---
% --- Archon physics formalization source end ---
